\documentclass[11pt,reqno]{amsart}
\usepackage[T1]{fontenc}
\usepackage{lmodern}
\usepackage[margin=1.15in]{geometry}
\setlength{\emergencystretch}{1.2em}
\usepackage{amsmath,amssymb,array,booktabs,longtable}
\usepackage[hidelinks]{hyperref}
\hypersetup{pdftitle={Off-centre Kaprekar cycles with positive basin density},pdfauthor={Rustem Galyavin}}
\newtheorem{theorem}{Theorem}
\newtheorem{lemma}[theorem]{Lemma}
\newtheorem{definition}[theorem]{Definition}
\theoremstyle{remark}
\newtheorem*{remark}{Remark}
\newcommand{\PP}{\mathbb P}
\newcommand{\ZZ}{\mathbb Z}
\newcommand{\lcm}{\operatorname{lcm}}
\newcommand{\sgn}{\operatorname{sgn}}
\newcommand{\file}[1]{\nolinkurl{#1}}
\title[Off-centre Kaprekar cycles]{Off-centre Kaprekar cycles with positive basin density}
\author{Rustem Galyavin}
\thanks{\href{mailto:rustem@galyavin.com}{rustem@galyavin.com}. ORCID \href{https://orcid.org/0009-0006-6875-9170}{0009-0006-6875-9170}.}
\date{20 November 2026}
\subjclass[2020]{11A63,37P99}
\keywords{Kaprekar map, basin density, digit differences, folded doubling}
\begin{document}
\raggedbottom
\begin{abstract}
Choose an ordered nonconstant word of fixed odd length \(n\) uniformly in
base \(B\), allowing leading zeros. At the centred profile the
normalized differences between the \(i\)th largest and \(i\)th smallest
digits, together with the endpoints 0 and 1, are equally spaced. At each of 47 explicitly listed lengths \(n\) there are
\(B_n\) and \(c_n>0\) such that, for every base \(B\ge B_n\), with
probability at least \(c_n\) the orbit of the word ends on a Kaprekar
cycle whose peak normalized distance from the centred profile exceeds a
fixed \(\delta_n>0\); for the inputs constructed in the proofs, every
state of the orbit stays at least that far. The
47 lengths are all those whose folded order (the least \(r\) with
\(2^r\equiv\pm1\pmod n\)) lies between three and thirteen and which admit
a gap profile, invariant under folded doubling, that satisfies certain
strict inequalities; each is certified by a finite computation, and
larger folded orders are not treated. At each of lengths 45 and 91, for each \(p\in\{24,132\}\), the
probability of reaching a cycle of least period \(p\) near the same
profile is bounded below by a positive constant for all sufficiently
large bases. At thirteen
digits two disjoint sets of inputs, each of probability bounded below
for all large bases divisible by seven, reach the same period-177 cycle
in those bases. Every phase of this last cycle has normalized radius at
most \(24/B\), which tends to zero. The proof
constructs full-rank integer families by inserting guarded translations
into finite reference trajectories and then realizes the resulting
orbits as actual subtractions with exact input weights.
\end{abstract}
\maketitle

\section{Introduction}

The Kaprekar map \(K_B\) sorts the digits of a word in decreasing and
increasing order and subtracts the second arrangement from the first
in base \(B\ge2\). Keep the word length \(n=2k+1\) fixed, pad with leading
zeros, and choose the original word uniformly from the \(B^n-B\)
nonconstant ordered words. Write \(m=k+1\). For sorted digits \(a_1\ge\cdots\ge a_n\),
set \(q_i=a_i-a_{n+1-i}\) for \(1\le i\le k\), with gap vector
\(g=(B-q_1,q_1-q_2,\ldots,q_k)\).
At the centred profile the normalized differences between the \(i\)th
largest and \(i\)th smallest digits, together with the endpoints 0 and 1,
are equally spaced; the vector of normalized differences is
\(((m-i)/m)_{i=1}^k\), the state with all gaps equal to \(B/m\), which we
call the centred state. The normalized radius of a state and the peak
normalized radius of the terminal cycle are
\[
 \mathcal R_B(q)=\max_{1\le i\le k}
        \left|\frac{q_i}{B}-\frac{m-i}{m}\right|,
 \qquad
 Z_B=\max_{q\ \mathrm{on\ the\ terminal\ cycle}}\mathcal R_B(q).
\]
The terminal cycle is the unique cycle eventually visited by the word.
The \emph{basin} of a cycle, or of a family of cycles, is the set of
original words whose orbit reaches it.
Let \(P_B\) be the least whole-word period of the terminal cycle and
\(\PP_B\) the original word probability. The lower basin density of a
cycle or family is the liminf of the probability of its basin on the
stated base domain; positive lower density means that this liminf is
positive.

At eleven digits, with probability tending to one through all bases,
an input reaches the paired region
\(\min_j g_j\ge\lceil(B+1)/7\rceil+64\) within \(4509B\) subtractions
\cite[Theorem~1]{G6}. For \(B\ge3534\), every phase of a terminal cycle
reached after such an entry has normalized radius at most \(420/B\)
\cite[Section~6.2]{G6}.
In particular \(Z_B\to0\) in probability there. The next two results
show that \(Z_B\) does not tend to zero in probability at the lengths
they cover. Their lower bounds are tiny and say nothing about the
typical input. The strict-length construction also keeps its certified
source sets outside a fixed neighbourhood of the centre throughout
their histories (Section~\ref{sec:lift} and the Remark after
Theorem~\ref{th:strict}). The separate thirteen-digit
result concerns a microscopic cycle and is compatible with macroscopic
centring.

\begin{theorem}[Fifteen and seventeen digits]\label{th:eight}
Let \(n=15\) or \(17\), and set \(j_1=1\) and
\(j_{i+1}=\min(2j_i,n-2j_i)\) for \(1\le i<4\). The relative
signature of a cycle is the lexicographically smallest tuple
\((0,g_{j_1}-g_0,\ldots,g_{j_4}-g_0)\) on the cycle.
Let \(E_B\) be the event that the terminal cycle has least period
eight, \(Z_B>1/200\), and signature \((0,-1,2,-1,0)\). Then
\[
 \liminf_{B\to\infty}\PP_B(E_B)
 \ \ge\ \frac{n!}{11^{(n+1)/2}\,100000^{(n-1)/2}}.
\]
\end{theorem}

\begin{theorem}[Strict lengths]\label{th:strict}
At each of the 47 strict lengths of folded order \(3\le r\le13\)
(the least \(r\) with \(2^r\equiv\pm1\pmod n\); Definition~\ref{def:strict}), the constants \(\delta_n,D_n\) of
Table~\ref{tab:lengths} satisfy
\[
  \liminf_{B\to\infty}\PP_B\{Z_B>\delta_n\}
      \ \ge\ \frac{n!}{(n+3)^{(n+1)/2}D_n^{(n-1)/2}}>0 .
\]
All integer bases are included. Explicitly, for each such \(n\) let
\(c_n\) be half the right-hand side; then \(c_n>0\) does not depend
on \(B\), and there is \(B_n\) with \(\PP_B\{Z_B>\delta_n\}\ge c_n\)
for every integer \(B\ge B_n\).
The 47 lengths are all the strict lengths of folded order 3 to 13
(Lemma~\ref{lem:compscreen}), and each is certified by a finite
reference record (Lemma~\ref{lem:compstrict}). Strict lengths of
folded order above 13 are not treated. The list includes nine digits
and the two lengths in Theorem~\ref{th:eight}.
\end{theorem}

\begin{remark}[Radius along the whole orbit]
The events in Theorems~\ref{th:eight}--\ref{th:coexist} concern the
peak radius \(Z_B\), and they may contain words outside the source
sets constructed in their proofs. For the constructed sets the proofs
give more. In Theorems~\ref{th:strict} and~\ref{th:coexist}, every
state of the orbit, from the original word through every phase of
the terminal cycle, has \(\mathcal R_B\ge3\delta_n/2\)
(Section~\ref{sec:lift}). In Theorem~\ref{th:eight} the bound on
\(g_0\) holds along the whole path, so every such state has
\(q_1/B-(m-1)/m>1/200\) (Section~\ref{sec:eight}). In particular
these terminal cycles stay outside the corresponding neighbourhood of
the centred state at every phase, not only at their peak.
\end{remark}

\begin{theorem}[Two periods at one profile]\label{th:coexist}
At \(n=45\) and \(91\), cycles near the same strict order-12 profile
have positive lower basin density for each of the least periods
\(p=24\) and \(132\). More precisely,
\[
 \liminf_{B\to\infty}\PP_B\{P_B=p,\ Z_B>\delta_n\}
  \ \ge\ \frac{n!}{(n+3)^{(n+1)/2}D_{n,p}^{(n-1)/2}},
\]
where
\[
\begin{array}{rccc}
 n & \delta_n & D_{n,24} & D_{n,132} \\
 \midrule
 45 & 1/598       & 87262111872  & 92890047744 \\
 91 & 119/285384  & 867599893776 & 923554563888
\end{array}
\]
and the liminf is over all integer bases.
\end{theorem}

\begin{theorem}[Thirteen digits in bases divisible by seven]\label{th:177}
As \(B\to\infty\) through \(B\equiv0\pmod 7\), there are two disjoint
sets of original words that reach the same least period-177 cycle.
The first set has lower limiting probability at least
\[
 \eta=\frac{3575}{1350641342287205909485672071168},
\]
and the second has limiting probability \(\mu>0\), given exactly by
\eqref{eq:mu}. In particular the cycle has lower basin density at
least \(\eta+\mu\) on this base domain. Every phase of this cycle has normalized radius at most \(24/B\).
\end{theorem}

\paragraph{Idea of the proofs.}
Every gap state represents at most \(O_n(B^{k+1})\) original words, so a
positive lower density requires \(\Omega_n(B^k)\) gap states
(Lemma~\ref{lem:route}). At a strict length, on the paired branches
with all secondary comparisons \(d_i\ge-1\), the \(r\) root differences
relative to the anchor \(g_0\) evolve autonomously. The complementary
gaps follow their permutation with a loss \(h\) at slot 3. On the
terminal zero-loss cycle their evolution is a pure permutation. Here
\(r\) is the length of the orbit of 1 under folded doubling, the root
coordinates are \(g_0\) and the \(r\) gaps indexed by that orbit, and
\(c=k-r\) is the number of complementary gaps. Root paths can run for a time
proportional to \(B\) before they settle into a cycle. A finite
reference path that retires its \(r\) root labels one at a time shows
that inserting up to \(L\) copies of suitable translating words sends
an \(L^r\)-point box of relative root sources to one zero-loss cycle
(Lemma~\ref{lem:retire}). With one anchor parameter and \(c-1\) free
complementary gaps, the last gap being fixed by the sum \(B\), this
gives \(L^{r+1+(c-1)}=L^k\) sources with \(L\) proportional to \(B\):
full rank. The route-count lemma then turns full rank into a uniform
positive lower bound. The leading real dynamics alone does not determine these basin weights.
The integer borrow corrections determine which terminal cycle a long
transient reaches.

\paragraph{Related work.}
Yamagami and Matsui used the order of 2 up to sign for ternary Kaprekar
loops \cite{YM}. For primitive proper loops of odd length \(n\) in base 3
the order already occurs modulo \(n-2\), as reviewed by Kay and
Downes-Ward \cite[Section~2]{KDO}, who also credit \cite{YM} before
their even-base Definition~7 \cite{KDE}.
A length-dependent modulus thus predates the present paper.
In even bases \(b\ge4\), Kay and Downes-Ward apply folded doubling
modulo \(b-1\) to the interior digit multiplicities of symmetric and
almost-symmetric full indices. The resulting periods divide the least
\(s>0\) with \(2^s\equiv\pm1\pmod{b-1}\)
\cite[Theorem~5(c), Definition~7 and Theorem~9]{KDE}. For odd bases
\(B>3\), Chen, Ono, Schwartz and Thakur classify the nonconstant
four-digit cycles by projective doubling on unordered pairs of nonzero
sign classes modulo \(B\). Each period divides the projective order
modulo \(B\) \cite[Theorem~1 and proof of Lemma~20]{COST}.
What is distinct here is the modulus \(n\) acting on gaps, the exact
borrow dynamics, and the growing-base word weights. Folded doubling
permutes gap coordinates, and the exact borrow corrections supply a
separate map on relative gaps. This distinction permits different
physical periods at a single folded order.

Devlin and Zeng give the fraction \(8/(5B^2)+8/(25u^2)\)
of all \(B^4\) four-digit inputs reaching the unique nonzero fixed
point at \(B=5u2^a\) with \(u>1\) odd and \(a\ge0\) \cite[Theorem~2]{DZ}.
Dahl enumerates individual attractor-basin weights at lengths five and
six in base ten \cite[Section~3.1 and Appendix~A.3]{Dahl}.
Cycle classifications and arrival laws in growing bases appear in
\cite{G1,G2,G3,G5}. The eleven-digit results give concentration at
the centre and arithmetic dependence of joint input and arrival-time
laws \cite{G6,G7,G8}. At nine digits the limit laws of
\cite[Theorem~2]{G4} already give, in each base class modulo 5, positive
limiting mass of inputs whose first cyclic word stays a fixed distance
from the centre; the nine-digit row of Theorem~\ref{th:strict} recovers
a positive lower bound there by the present method, with an explicit
but much smaller constant, and adds the radius bound along the whole
orbit of the constructed set. New here are: the all-large-base lower
bounds at the 46 other listed lengths; the period-eight family at 15
and 17 digits; the two period classes 24 and 132 at one profile at 45
and 91 digits, with the physical period fixed separately by the root
and the complement; and the microscopic period-177 cycle at thirteen digits
with two disjoint entrance sets. None of this requires a limiting
distribution on the entire population. We prove nothing about the
remaining population at these lengths and, apart from the
thirteen-digit result, no positive lower density at lengths outside
the list.

\section{Counting and exact subtraction}\label{sec:framework}

The integer gap vector has \(m\) nonnegative entries and sum \(B\), with
\(g_0=B-q_1\ge1\) and
\(q_i=\sum_{j=i}^k g_j\). An original word determines its gaps uniquely.
Conversely, the numerical subtraction is
\(\sum_{i=1}^kq_i(B^{n-i}-B^{i-1})\), so the gaps determine the next
word. Projection to gaps and reconstruction of the next word give
both divisibilities between the least periods on cycles. Thus
a certified physical gap cycle has the same least period as its
whole-word cycle.

For positive integer gaps the number of original ordered words over
a gap vector, with all digits distinct, is exactly
\begin{equation}\label{eq:count}
                         n!\,g_0\prod_{i=1}^k(g_i-1).
\end{equation}
Indeed, choose the smallest digit from \(0,\ldots,g_0-1\).
The total digit spread is \(B-g_0\), so every allowed smallest
digit keeps the largest at most \(B-1\). Each \(g_i\), \(i>0\),
is the sum of two sorted spacings placed symmetrically about
the middle digit. Choose its positive split
independently, then order the \(n\) different digits arbitrarily.
Allowing zero spacings gives the upper bound
\(n!g_0\prod_{i=1}^k(g_i+1)\) for all words over the state.
Words with a zero spacing contribute at most \(O_n(B^k)\) per
gap state by fixing one spacing to zero before counting the others.

\begin{lemma}[Route count]\label{lem:route}
Let a source chart at fixed length and on a specified base domain
be an affine integer parametrization with \(s\) independent directions.
Suppose each parameter ranges over an interval of length
\(\Theta(B)\), and that for a fixed \(c>0\) every initial gap vector
of the chart satisfies \(g_0\ge cB\) and \(g_i-1\ge cB\) for
\(1\le i\le k\). Suppose that every chart source follows a certified
physical itinerary to a specified event. The chart has probability
\(\Theta(B^{s-k})\) in the original ensemble. In particular at rank
\(s=k\) it gives a uniform positive lower bound on the probability
of that event.
\end{lemma}
\begin{proof}
Injectivity gives \(\Theta(B^s)\) gap states. Each has
\(\Theta(B^{k+1})\) original words by \eqref{eq:count} and the
upper count. Divide by \(B^{2k+1}-B\). The lower bound uses only the
exact count \eqref{eq:count} and the hypothesis on the gaps; the upper
count is needed only for the matching upper order.
\end{proof}

An itinerary specifies successive affine branches, their fixed
connecting words, and the affine counts of repeated words. A
\emph{guard} is an inequality specifying the branch's borrowing or
sorting comparisons. Once a branch and tie convention are fixed,
these comparisons are affine functions of the input. For a translating
word with reverse displacement \(v\) and target \(y\), the state with
\(j\) copies remaining is \(y+jv\). Verifying the comparisons at
\(j=1\) and \(j=c\) suffices for all copies \(1\le j\le c\).
The output identities at those endpoints likewise imply the
intermediate identities by affinity. When the source parameters
range over boxes, check the minimum of each comparison exactly by
choosing the appropriate interval endpoint for each coefficient.
This proves the itinerary by induction for every lattice point of
the chart; following a sample trajectory by literal subtraction serves
as an additional check.

\subsection{When the borrowed pairs can be sorted separately}

If \(q_k>0\), the primitive borrowed digits are the maximal
digit \(B-1\), from the middle column, and the \(k\) pairs
\begin{equation}\label{eq:pairs}
  \left(q_i-\mathbf1_{i=k},\quad
               B-q_{k+1-i}-\mathbf1_{i<k}\right),\qquad 1\le i\le k .
\end{equation}
This follows by propagating the borrow from the least significant
position through the middle digit to the leading position. The state is
\emph{paired} if this maximal digit is largest and both digits of pair
\(i\) are at least both digits of pair \(i+1\).
A sufficient strict condition is
\begin{equation}\label{eq:physical}
                     (m+1)\min_j g_j>B+1.
\end{equation}
Indeed, write \((\lambda_i,\lambda'_i)\) for pair \(i\) in \eqref{eq:pairs}.
For \(1\le i<k\) the four comparisons between pairs \(i\) and \(i+1\) are
\[
\begin{aligned}
 \lambda_i-\lambda_{i+1}&=g_i+\mathbf1_{i+1=k}, &
 \lambda'_i-\lambda'_{i+1}&=g_{k-i}-\mathbf1_{i+1=k},\\
 \lambda_i-\lambda'_{i+1}&=q_i+q_{k-i}-B+\mathbf1_{i+1<k}, &
 \lambda'_i-\lambda_{i+1}&=(B-q_{k+1-i})+(B-q_{i+1})-B-\mathbf1_{i+1<k}.
\end{aligned}
\]
The same-column comparisons are at least \(\min_j g_j-1\).
In the cross comparisons, \(q_i+q_{k-i}=\sum_{j\ge i}g_j+\sum_{j\ge k-i}g_j\)
is a sum of two suffixes of \(g\), and
\((B-q_{k+1-i})+(B-q_{i+1})=\sum_{j\le k-i}g_j+\sum_{j\le i}g_j\) is a sum
of two prefixes; in each case the two sums together have \(m+1\) terms.
So each cross comparison is at least \((m+1)\min_jg_j-B-1\).
Positive gaps also keep \(B-1\) at the head of the
ordered list. Thus sorting within each pair and taking symmetric
differences gives the actual next gap vector. The weak inequality
\((m+1)\min_j g_j\ge B+1\) also suffices, with ties between successive
levels allowed.

Let \(P\) be the permutation sending coordinate \(j>0\) to
\(f_n(j)=\min(2j,n-2j)\) and fixing \(0\). Folded doubling is
multiplication by 2 on nonzero sign classes modulo \(n\), in the
representatives \(1,\ldots,k\). Put
\[
 d_i(g)=\sum_{j=m-i}^{m-1}g_j-\sum_{j=0}^{i-1}g_j,
 \qquad 1\le i\le\lfloor m/2\rfloor .
\]
The two reflected comparisons in \eqref{eq:pairs} have the same
\(d_i\); the interior differences have the offset \(+1\).
We use the part of the paired region with
\(d_i\ge -1\) for \(2\le i\le\lfloor m/2\rfloor\).
On it, write \(z=d_1=g_k-g_0\), \(e=\sgn z\), and \(h=(-1-z)_+\).
The exact update is
\begin{equation}\label{eq:borrowed}
      g^+=Pg+e(1,-2,1,0,\ldots)+h(-1,1,1,-1,0,\ldots).
\end{equation}
Here and below the root map is used at orders at least three.
For a direct derivation, first treat the last primitive pair
provisionally as \((q_k,B-q_1-1)\). Sort within each pair into its
maximum and minimum, retaining the given levels. The provisional gaps
at indices \(2i,2i+1\) are given by the following formula, with the
plus sign at \(2i\):
\begin{equation}\label{eq:maxmin}
 \frac{g_i+g_{k-i}\ \pm\bigl( |d_i+1|-|d_{i+1}+1|\bigr)}2 ,
 \qquad 0\le i<\lfloor m/2\rfloor ,
\end{equation}
where \(d_0=0\); append \(g_{k/2}\) if \(m\) is odd. The unequal
borrow adjustment on the first and last primitive pair changes
the first three gaps by \(e,-2e,e\).
For example, the first gap of the displayed list is
\(\min(g_0,g_k+1)\), and the exact first gap is
\(\min(g_k,g_0+1)\). Their difference is \(\sgn(g_k-g_0)\) on
integer gaps. The next two differences follow by the same
max/min comparison of the end pairs. Substituting \(d_i+1\)
for all interior absolute values except \( |d_1+1|=d_1+1+2h\)
gives \eqref{eq:borrowed}. The corresponding leading operation
omits the integer borrow terms and has the majorization property
of \cite[Theorem~1]{G4}. The loss \(h\) here measures a transfer
between coordinates in the exact integer map.

\subsection{Strict lengths and the root map}

Let \(J=\{0\}\cup C\), where \(C\) is the orbit of 1 under \(f_n\),
listed in order of successive folded doubling starting at 1 and
ending at \(k\). It ends there because \(f_n(k)=1\). Its length
\(r=|C|\), the \emph{folded order}, is the least \(r>0\) with
\(2^r\equiv\pm1\pmod n\).

\begin{definition}\label{def:strict}
A length \(n\) is \emph{strict} if there is a real profile \(\pi\)
on the gap simplex, invariant under \(P\) and constant on \(J\),
such that
\(d_i(\pi)>0\) for all \(2\le i\le\lfloor m/2\rfloor\).
A \emph{full-orbit length} is one with \(J=\{0,\ldots,k\}\).
A \emph{forced-tie length} in the range considered here has
\(J\ne\{0,\ldots,k\}\), but some secondary comparison \(d_i\)
vanishes identically on the invariant profile space constant on \(J\).
\end{definition}

Every \(d_i\) vanishes on the constant profile. Folded order \(r\ge3\)
forces \(n\ge7\), so \(\lfloor m/2\rfloor\ge2\) and the condition in
the definition is not vacuous; hence a strict profile is not constant,
and \(A>0\) in Section~\ref{sec:lift}.

At a strict length, the index 3 belongs to the complement of \(J\): its
predecessor \(k-1\) under \(P\) would otherwise make \(d_2=0\).
For the differences \(x_i=g_{j_i}-g_0\), the root (that is,
the coordinates in \(J\) relative to their anchor \(g_0\)) has
an autonomous integer map
\begin{equation}\label{eq:root}
   T(x)=(z-3e+2h,\ x_1+2h,\ x_2-e+h,\ldots,x_{r-1}-e+h),
   \qquad z=x_r .
\end{equation}
The first coordinates in \(C\) are 1 and 2, so subtracting the
updated anchor from their updates in \eqref{eq:borrowed} gives
the displayed formula. The root sum in absolute coordinates
increases by \(h\); the complement follows its permutation with
\(-h\) at index 3. A zero-loss root cycle is one for which
the sum of \(h\) over a full period is zero. Since all losses are
nonnegative, its absolute root mean is then constant. Root
recurrence with strictly positive loss does not itself give a
physical cycle.

\section{A common period-eight basin at lengths 15 and 17}\label{sec:eight}

The first example uses the root orbit \(1,2,4,k\). The two profiles
take value \(\alpha\) on \(J\) and \(\beta\) off it:
\[
 \begin{array}{rcc}
 n & \alpha & \beta\\
 \midrule
 15 & 47/400 & 47/400+1/50 \\
 17 & 14/135 & 14/135+1/60
 \end{array}
\]
The same four-parameter root chart works at both lengths.
For \(0\le a,b,c,d<L=\lfloor B/100000\rfloor\), put
\(v=5a+12b\), and start at
\[
       x=(5a, 8+v, 49+27(d+1)+26(c+1)+v, 31+26(c+1)+v).
\]
Its determinant has absolute value \(42120\). The five sign words of
Appendix~\ref{app:eight}, in order, are used \(a,b,1,c+1,d+1\)
times. They lead to \(x=(0,0,18,2)\). A 35-step prefix there
leads to the zero-loss eight-cycle with listed signature. The
root spread \(\max(0,x)-\min(0,x)\) is at most \(179L\) along
the whole path. The cumulative loss is
\(4b+18(c+1)+13(d+1)+22\le53L\).

Choose the anchor among \(L\) integers at the floor of
\(B\alpha-\sum_i x_i/5\), and each complementary gap but one
among \(L\) integers at the floor of \(B\beta\).
The root mean stays within \((1+53/5)L\) of \(B\alpha\);
together with spread this gives root error at most \(191L\).
The dependent complementary gap begins within \(8L\)
of \(B\beta\), by the base-sum constraint, and acquires at most
\(53L\) of additional error. In the
two profiles every secondary comparison has margin
at least \(\beta-\alpha\) and uses at most eight terms.
The inequalities
\[
\begin{gathered}
(m+1)(\alpha-191/100000)>1,\qquad
8(191/100000)<\beta-\alpha,\\
1/m-\alpha-191/100000>1/200
\end{gathered}
\]
preserve pairing, secondary guards and the radius (using \(q_1/B=1-g_0/B\)).
The complementary permutation orders are two and four,
respectively, so the physical least period is eight.
There are \(L^k\) gap states, with \(g_0,g_i-1\ge B/11\)
for all sufficiently large bases. Equation~\eqref{eq:count}
proves Theorem~\ref{th:eight}.

\section{Retirement and lifting at all strict lengths}\label{sec:lift}

Let \(S\) rotate the relative coordinates of \eqref{eq:root}, moving
the last coordinate to the first. Label coordinates by their original
slots, so that at the \(t\)th step each label occupies its slot under
\(S^t\). Its endpoint read is the step at which it supplies \(z\).
A label is called \emph{active} until its first nonpositive read;
that first read is its \emph{retirement}. In a projected return word,
ignore active values and prescribe the positive branch at their reads.
A return of length divisible by \(r\) restores the other values to
their initial slots and values. It translates the active coordinates
by the accumulated kicks, provided their reads stay positive.

\begin{lemma}[Retirement]\label{lem:retire}
Suppose a finite reference path to a zero-loss root cycle retires
all \(r\) labels one by one. In each stage suppose there is a projected
return word in the reference path, ending at or before the next
retirement, which strictly decreases the next retiring coordinate.
Let \(u_i\) be the reverse displacement of that word, transported to
the original slots. The columns \(u_i\), in stage order, form a
triangular array when rows are put in retirement order.
Its absolute determinant is the product of the positive retiring
drifts. There are \(L^r\) distinct guarded root sources, all entering
the same zero-loss cycle, with time, spread and loss \(O_r(L)\)
up to and through the cycle.
\end{lemma}
\begin{proof}
A return restores every previously retired label, so \(u_i\) vanishes
there. Its entry for the next retirement is the positive reverse drift.
These are the triangular array and its diagonal. Start from the reference
source plus \(\sum_i a_i u_i\).
Take \(a_i\) in \(L\) consecutive integers starting at \(K_iL\).
Choose the fixed integers \(K_i\ge0\) in stage order. For the label
retiring at stage \(i\), choose \(K_i\) large enough that its positive
diagonal term dominates the largest negative sum of all earlier columns.
Later columns vanish on that label. While earlier copies are being
consumed their remaining negative contributions are no larger. During
stage \(i\) the earlier copies are gone. Thus the perturbation at a
positive read is nonnegative while any copies remain, and at a
nonpositive read the perturbation is zero. It follows from
\eqref{eq:root} that the perturbation only rotates and the loss follows
the prescribed word. Indeed, on the branch \(z\ge1\) the map is \(S\)
plus a constant, on \(z=0\) it is \(S\), and on \(z\le-1\) the read
value \(z\) enters every output while the other coordinates are only
rotated. Hence, if \(p_r\ge0\) at a positive read of \(x\) and
\(p_r=0\) at a nonpositive read, then
\begin{equation}\label{eq:rotate}
 T(x+p)=T(x)+Sp,\quad\text{with the same loss } h .
\end{equation}
Write the reference path as \(x(t)\) and the stage-\(i\) return as the
interval \([s_i,t_i)\). Its length is divisible by \(r\), so
\(S^{t_i}=S^{s_i}\) and \(x(t_i)=x(s_i)-S^{s_i}u_i\). Iterating
\eqref{eq:rotate} along the interval sends a state
\(x(s_i)+S^{s_i}p\), with \(p\) in original label coordinates, to
\(x(t_i)+S^{t_i}p=x(s_i)+S^{s_i}(p-u_i)\). Thus an extra complete
return at the start of stage \(i\) subtracts one \(u_i\) in original
label coordinates and comes back to the same reference phase. From
\(x(0)+\sum_ia_iu_i\), follow the reference path to \(s_1\), traverse
\([s_1,t_1)\) \(a_1\) extra times, continue to \(s_2\), and so on
through stage \(r\). After the last extra return the perturbation is
zero, and from that phase on the path coincides with the reference
path. Hence the terminal state and the cycle reached are those of
the reference path and do not depend on the copy counts \(a_i\).
Invertibility gives \(L^r\) sources. The clock and cumulative loss grow by the lengths
and losses of the inserted words; spread is bounded by the reference
spread plus the sum of spreads of all remaining rays. These are all
linear in \(L\).
\end{proof}

For the finite certificates it is enough to check each inserted word
against all its earlier reads. In a reference orbit \(x(t)\) an interval
\([s_i,t_i)\) supplies
\(u_i=S^{-s_i}(x(s_i)-x(t_i))\).
If \(t_i-s_i\) is a multiple of \(r\), every affected label is read
positively up to \(t_i\), and the array of rays is nonnegative and
full rank, the proof above works with \(K_i=0\). One can therefore
verify a certificate with finite exact replays and then vary the
copy counts over arbitrarily long intervals. Appendix~\ref{app:strict}
gives the finite predicates and budget definitions.

\subsection{Realizing the root paths in every large base}

A \emph{physical lift} assigns the anchor and complementary gaps
to a relative root chart and proves that the resulting full-gap path
is actual digit subtraction. Let \(\pi\) be a strict rational profile,
mixed toward the constant profile so that
\[
 \Delta=\min_j\pi_j-\frac1{m+1}>0,\quad
 \gamma=\min_{i\ge2}d_i(\pi)>0,\quad
 A=\max_{1\le i\le k}\left|\sum_{j=i}^k\pi_j-\frac{m-i}{m}\right|>0.
\]
Suppose an order-\(r\) certificate gives \(L^r\) root sources with
spread and cumulative loss each at most \(CL\).
There are \(c=m-r-1\) complementary gaps. For each root choose the
anchor among \(L\) integers starting at
\(\lfloor B\pi_0-\sum_i x_i/(r+1)\rfloor\).
Choose \(c-1\) complementary gaps independently in intervals of size
\(L\) starting at their profile floors. The remaining gap enforces
the sum \(B\). It starts with error at most \((m-1)L\).
Indeed, the root sum has initial error at most \((r+1)L\), and the
free complementary gaps account for at most \((c-1)L\).
The root mean changes only by the total loss divided by \(r+1\).
The complement follows a permutation with accumulated error at
most the total loss. Every formal gap along the path therefore
differs from \(B\pi\) by at most \(EL\) with \(E=m+2C+1\).
Choose
\begin{equation}\label{eq:D}
       D=\left\lceil\frac{4E}{\min(\Delta,\gamma/m,A/m)}\right\rceil,
                       \qquad L=\lfloor B/D\rfloor.
\end{equation}
The error is at most a quarter of each indicated margin in
normalized gap units. In particular \eqref{eq:physical} and the
secondary guards hold for all sufficiently large bases, at every step
by induction. The radius is at least \(3A/4>A/2\) on the whole
constructed route, including the completed cycle. The constant
root mean and the complementary permutation realize the recurrent
zero-loss root as a physical cycle. No congruence obstruction
arises because a free complementary gap enforces the integer base sum.

The independent choices give \(L^{r+c}=L^k\) gap states. At all
sufficiently large bases each of \(g_0,g_i-1\) exceeds
\(B/[2(m+1)]\). Equation~\eqref{eq:count} gives the claimed lower
bound with \(\delta_n=A/2\). The all-base lift explains why
periodic points of a leading real map are not enough by themselves
to decide the original population question: transient paths
of length proportional to \(B\) give these families their full
source rank.

\subsection{The finite length screen}

For a given order \(r\), every candidate \(n\) divides \(2^r-1\)
or \(2^r+1\). After testing for exact order, assign arbitrary
weights on the complementary permutation orbits and set the
weight on \(J\) to zero. The folded-fixed vectors constant on \(J\)
are spanned by such vectors and the all-ones vector, and all
\(d_i\) annihilate the latter. For the 47 strict lengths a single
integer witness at each length gives every secondary comparison
positive. For each other non-full length in the range, at least
one row annihilates every complementary indicator, proving a
forced tie. The witnesses and profile formula are in
Appendix~\ref{app:strict}; Table~\ref{tab:lengths} reports all bounds.

The five full-orbit lengths are \(7,11,13,19,23\), of orders
\(3,5,6,9,11\). Their invariant profile in this construction is
the constant one. The seven- and eleven-digit centring results
are in \cite{G5,G6,G8}, and Theorem~\ref{th:177} gives the separate
thirteen-digit result. The screen gives no population conclusion
at 19 or 23 digits. There are off-centre fixed points on a sequence of
bases at 23 digits, as shown in Section~\ref{sec:remarks}.
The four forced-tie lengths are \(25,27,35,39\), with vanishing
secondary rows \(\{2,6\},\{4\},\{2\},\{4\}\), respectively.
At 25 digits there is an all-large-base family of off-centre fixed
points and a feeder of mass \(\Theta(B^{-10})\) described in
Section~\ref{sec:remarks}. Positive lower density at these four
lengths remains open in this analysis.

% Table of the 47 strict lengths (inlined from lengths.tex)
\begin{longtable}{@{}rrrr@{}}
\caption{The 47 strict lengths and their lower basin bounds. Each row gives
\(\delta_n\) and \(D_n\), with mass bound
\(n!/((n+3)^{(n+1)/2}D_n^{(n-1)/2})\).}
\label{tab:lengths}\\
\toprule
Length \(n\) & Folded order \(r\) & Radius threshold \(\delta_n\) & Mass divisor \(D_n\)\\
\midrule
\endfirsthead
\caption[]{The 47 strict lengths (continued).}\\
\toprule
Length \(n\) & Folded order \(r\) & Radius threshold \(\delta_n\) & Mass divisor \(D_n\)\\
\midrule
\endhead
\midrule
\multicolumn{4}{r@{}}{\footnotesize Continued on the next page}\\
\endfoot
\bottomrule
\endlastfoot
9 & 3 & 1/200 & 7708000 \\
15 & 4 & 1/256 & 7736832 \\
17 & 4 & 1/270 & 10886400 \\
21 & 6 & 13/5808 & 246700608 \\
31 & 5 & 1/512 & 66533376 \\
33 & 5 & 19/10404 & 79569792 \\
41 & 10 & 47/38808 & 9997872192 \\
43 & 7 & 109/133584 & 717880416 \\
45 & 12 & 1/598 & 87262111872 \\
51 & 8 & 2/4563 & 7997661360 \\
57 & 9 & 11/33640 & 18105855360 \\
63 & 6 & 97/135168 & 5747073024 \\
65 & 6 & 1/1452 & 6297382080 \\
73 & 9 & 61/104044 & 18669655360 \\
85 & 8 & 67/162712 & 23776779136 \\
89 & 11 & 7/20700 & 214723418400 \\
91 & 12 & 119/285384 & 867599893776 \\
93 & 10 & 17/53016 & 109287606528 \\
105 & 12 & 25/101124 & 1844579423232 \\
117 & 12 & 33/139240 & 2539854561600 \\
127 & 7 & 31/133120 & 17258741760 \\
129 & 7 & 127/557700 & 18078403200 \\
171 & 9 & 173/1286904 & 231174286944 \\
195 & 12 & 373/3803184 & 11562508454112 \\
205 & 10 & 431/4413344 & 1137707123072 \\
241 & 12 & 15/162382 & 21722102549312 \\
255 & 8 & 205/2818048 & 619130781696 \\
257 & 8 & 69/961480 & 633734543520 \\
273 & 12 & 389/5180244 & 31498950224448 \\
315 & 12 & 727/15877104 & 48271762621152 \\
341 & 10 & 839/20117808 & 5188865510592 \\
455 & 12 & 1207/47617344 & 144779486850432 \\
511 & 9 & 731/33685504 & 6074035339264 \\
513 & 9 & 1469/68162568 & 6145537130880 \\
585 & 12 & 1675/100958424 & 306974991681792 \\
683 & 11 & 37/2971752 & 92574270913824 \\
819 & 12 & 2441/276356400 & 840356148088800 \\
1023 & 10 & 1127/179306496 & 139109207113728 \\
1025 & 10 & 3389/541075464 & 139926443594112 \\
1365 & 12 & 4619/1276313904 & 3881765791858752 \\
2047 & 11 & 1917/1074790400 & 2485955801907200 \\
2049 & 11 & 853/479085000 & 2493252240660000 \\
2731 & 13 & 1311/1275380926 & 9129135856118368 \\
4095 & 12 & 5713/11458838528 & 104646169866534912 \\
4097 & 12 & 17149/34426888200 & 104799578907384000 \\
8191 & 13 & 18937/137472507904 & 247507152900456448 \\
8193 & 13 & 37885/275146424328 & 247689012351460224 \\
\end{longtable}


\section{Periods 24 and 132 at lengths 45 and 91}\label{sec:coexist}

The order-12 reference route in Appendix~\ref{app:strict}
reaches a zero-loss root cycle of least period 24. Raising
the initial label in slot 7 by \(12592\) gives another reference
route with the same insertion rays and a least period-11
zero-loss root cycle. The finite replay verifies the new
suffix, the guards and the budgets of both entire root paths.
For each of the two word lengths the physical lift uses the
same strict profile for these two references.

At length 45 the complementary permutation has orbit lengths
\(4,3,2,1\); at length 91 they are \(12,6,12,3\).
Both permutations thus have order 12. At fixed original root
and anchor, the complementary terminal vector is an affine
translation of a permutation of the initial one. The \(c-1\)
free integers range independently across intervals of length \(L\).
An equality between any two terminal entries cuts out at most
\(O_n(L^{c-2})\) choices, since a difference of two entries
is nonconstant on the fixed-sum hyperplane. After excluding this
finite union, the complementary vector has least period 12.
The joint least periods are respectively
\(\lcm(24,12)=24\) and \(\lcm(11,12)=132\).
Discarding this lower-order number of sources leaves the limiting
counts in Theorem~\ref{th:coexist} unchanged.

\section{Two routes to a microscopic cycle at thirteen digits}\label{sec:13}

Write \(X_i=7q_i-(7-i)B\) and
\(H(X)=(-X_1,X_1-X_2,\ldots,X_5-X_6,X_6)\).
All \(X_i\) are multiples of seven in the base class in question,
and \(g_i=(B+H_i(X))/7\). Put \(R=-\min_iH_i(X)\).
The sufficient paired guard is \(B\ge8R+7\), by
\eqref{eq:physical}. The certified cycle has \(R\le28\).
Since \(\sum H_i=0\), any partial sum of \(H\) has absolute
value at most \(6R\); hence \(\mathcal R_B(q)\le24/B\) there.

The first source, in coordinates \(X/7\), is an integer box
on six independent rays with determinant \(84015792\).
Its six translating blocks and fixed bridges collapse the entire
box to one cycle phase. For parameters from 0 through
\(L=\lfloor B/90608\rfloor\), the paired bound along the route is
\(R\le22834+5663L\). At \(B\ge365360\)
this implies \(B\ge8R+7\) and \(g_j\ge B/8\).
Then \(g_i-1\ge B/9\) for \(i>0\), and
\eqref{eq:count} gives
\[
 \PP_B(\text{first source})\ge
 \frac{13!(B/8)(B/9)^6(1+\lfloor B/90608\rfloor)^6}{B^{13}-B}
 \ \ge\ \eta.
\]

The second source uses an affine box on a finer integer progression
in its repeat counts, while ranging over every sufficiently large
\(B\equiv0\pmod7\). The initial \(g_0<B/8\), so its original
words are disjoint from the first source. Five repeat counts absorb
directions as the itinerary passes sorting boundaries. The final
count cancels all remaining directions, including the scaling
parameter, at the fixed gap-offset state
\(H=(-1988,1484,1491,-1974,315,322,350)\).
The continuation from this state has preperiod 437 and the
same least period-177 cycle, verified by comparing the complete
cyclic state sets after rescaling \(X/7\) to \(X\).
The affine guard checks include the first and last copies,
the deciding comparisons at the five exits, and the first hit of
\(R\le2000\) from this box. Appendix~\ref{app:13} specifies the
two certificates.

For the second chart let \(B_0=140000\), \(S_0\) be its reference
source in \(X\), \(D_i\) its six direction columns, \(w_i\) the
half-widths of its normalized box and
\(M=63\lfloor(B-B_0)/(63B_0)\rfloor\).
The source is \((1+M)S_0+\sum_i z_iD_i\) for
\(|z_i|\le\lfloor Mw_i\rfloor\).
Since \(M/B\to1/B_0\) and the parametrization is injective,
\eqref{eq:count} and a Riemann sum yield
\begin{equation}\label{eq:mu}
 \mu=\frac{13!}{B_0^6}
 \int_{\prod_i[-w_i,w_i]}
      \prod_{j=0}^6\frac{B_0+H_j(S_0+\sum_i v_iD_i)}{7B_0}
       \,dv_1\cdots dv_6 .
\end{equation}
This uses the lattice of the \(z_i\) directly and requires no
unweighted approximation to the physical input measure.
The tied-digit correction is lower order by the bound following
\eqref{eq:count}. The constant \(\mu\) has the exact value given in
Appendix~\ref{app:13}, approximately \(1.11232\times10^{-21}\).
The nonnegative normalized radius of the first cyclic word tends
to zero on both sets, irrespective of the earlier route.

\section{Remarks and open questions}\label{sec:remarks}

\begin{remark}[Size of the bounds]
The quantitative lower bounds come from boxes small enough to
certify every comparison and every copy of each return. Many are
astronomically small. They give no estimate of the actual fraction
of words reaching the off-centre part of the terminal set.
The certified strict routes also give no asymptotic frequency
as the word length itself grows.
\end{remark}

\begin{remark}[A cube model for the root recurrence]
Let \(z_r(L)\) be the share of the uniform integer root cube
\([0,L)^r\cap\ZZ^r\) that ends at a zero-loss root cycle.
A budgeted version restricts to paths with clock, total loss
and spread each bounded by \(CL\) for a fixed \(C\).
Write \(\zeta(r)=\lim_{L\to\infty}z_r(L)\) when the limit exists.
The parity conjecture asks for positive lower density, with fixed
linear budgets allowed to depend on \(r\), at each fixed odd
order, and for \(\limsup_L z_r(L)\to0\) along even \(r\to\infty\).
Existence of the cube limits is a separate question.
An even order can have positive lower density while the decay
along the sequence of even orders still holds. A finite test of
36 prespecified comparisons confirmed 21, left 15 unresolved
and contradicted none. At order 8 and \(L=10^6\), 4305 of
8192 sampled sources had zero loss around their terminal cycle.
These are finite-scale results for uniformly sampled roots.
The original-word measure and the distribution of physical
entrances require a separate calculation.
\end{remark}

\begin{remark}[Two certified cohorts at order 8]
A component here is a cycle of the map obtained by ignoring
the last active label and forcing its positive branch.
Two projected components of periods 1712 and 368 have
descents 369 and 66 per complete return of that map.
In the respective certified subsets, the zero-loss terminal
phase fractions are \(194/369\) and \(39/66\).
The limiting uniform-cube weights of the subsets, including
all eight disjoint cyclic rotations of each retirement order,
are about \(5.24\cdot10^{-17}\) and \(3.21\cdot10^{-17}\).
Their exact volume formulas and scope are in
Appendix~\ref{app:cube}. At 51 digits each of two explicit
entrance families lifts to a positive original-word liminf
at radius above \(2/4563\). The lower bounds are
\(51!/(54^{26}D^{25})\) with \(D=7997661360\) and
\(50745451536\), respectively. These last bounds are obtained
by the guarded lift, with the original-word multiplicities
from \eqref{eq:count}.
\end{remark}

\begin{remark}[Unresolved access]
Can full-rank zero-loss basins be constructed by a common argument
for folded orders above 13? A downward projected drift identifies
an opportunity to retire a label. Its next component, and ultimately
its terminal loss, still depend on the retirement state.
For thirteen digits, original access and terminal selection at
the other base classes remain open in this paper.
\end{remark}

\begin{remark}[Two excluded lengths with macroscopic fixed points]
At 23 digits, take \(B=71a\) with \(a\ge7\) odd and
\(12q_i=(12-i)B+X_i\), where
\[
 X=(-a,-2a,-3a,-4a+12,a-18,-18,-a-6,-2a+6,-3a+18,2a-12,a).
\]
The gaps are integral and have minimum \((B-5a+30)/12\), so
\eqref{eq:physical} holds. The six comparisons \(d_i\) are
\(0,-1,(-a+3)/2,(-a+3)/2,-2,-3\).
Substituting in the paired max/min formula \eqref{eq:maxmin}, whose
end-pair correction vanishes here because \(d_1=0\), gives the same
eleven differences. In
particular the whole-word fixed point has
\(\mathcal R_B\ge a/(12B)=1/852\). This supplies a fixed point
on a sequence of bases and gives no estimate of its basin.

At 25 digits, let \(\tau\) be an integer in
\([B/1000,B/900]\) with \(2\tau\equiv B\pmod{13}\) and set
\(u=(B-2\tau)/13\). The profile with gaps \(u\) except for
\(g_5=g_{10}=u+\tau\) is fixed by \eqref{eq:borrowed}.
Indeed its comparisons are \(0,0,\tau,\tau,\tau,0\).
For an integer \(j\) between 0 and \(B/10^6\), decrease the
anchor by \(10j\) and increase each of the other ten root
gaps by \(j\). In a ten-step round with \(j=1\), the triples
\((d_1,d_2,d_6)\) are
\[
\begin{gathered}
(11,11,11),(10,12,11),(9,11,11),(8,10,13),(7,9,11),\\
(6,8,9),(5,8,7),(4,5,3),(3,4,4),(3,4,4).
\end{gathered}
\]
They follow by iterating \(Pg+(1,-2,1,0,\ldots)\).
Each additional unit of \(j\) adds eleven to the three
comparisons, since the perturbation is constant on the
non-anchor root orbit. The other comparisons stay at least
\(\tau-13(10j+20)>0\), and the pairing condition holds for
large \(B\) since \(\min g_i\ge u-10j-20\).
One round adds ten to the anchor and subtracts one from each
non-anchor root gap, so it decreases \(j\) by one.
These sources have two independent integer parameters and
reach fixed points of radius at least \(2\tau/(13B)>1/7000\).
Lemma~\ref{lem:route} gives their original mass
\(\Theta(B^{2-12})\).
\end{remark}

\section{Computations and verification}\label{sec:verification}

The file \file{ancillary/reproduce.sh} verifies the pinned inputs
and runs the exact checks in a fresh temporary directory.
From \file{ancillary/}, run \texttt{bash reproduce.sh}.
The proof computations use integer dynamics, rational Gaussian
elimination, guard inequalities and rational integration. They
use no external solver. The tested language and compiler
versions are recorded with the ancillary replay.
An independent reconstruction, with separately written programs
included in the ancillary files, recomputes the constants with its own
arithmetic and includes literal whole-word controls. Both implementations use
the pinned finite descriptions of the reference paths. The independent
checks cover these certificates and their quantitative consequences.
A complete classification of cycles and a proof of the parity
conjecture remain outside the verification scope.

\begin{center}
\small
\begin{tabular}{@{}>{\raggedright\arraybackslash}p{0.21\linewidth}>{\raggedright\arraybackslash}p{0.34\linewidth}>{\raggedright\arraybackslash}p{0.37\linewidth}@{}}
\toprule
Certificate & Exact primary check & Independent coverage \\
\midrule
Period eight
& Affine sign-word guards, rank, loss, spread, terminal period
and physical margins.
& Reconstructs the affine words and both lifts,
including the signature and literal controls.\\[3pt]
The 47 lengths
& Exact divisor screen, rational profiles, root insertions,
linear budgets and retirement determinants.
& Enumerates all candidates again, reconstructs the
root rays and all 47 pairs \((\delta_n,D_n)\), and checks
representative paths by literal subtraction.\\[3pt]
Order-12 coexistence
& Reconstruct the period-24 certificate from its itinerary and
verify the shifted period-11 certificate.
& Reconstructs the shifted lift, verifies generic
complement periods and checks both physical periods.\\[3pt]
Thirteen digits
& Both affine routes, paired guards, source weights,
and the 177-cycle, including agreement of the two terminal cycles.
& Reconstructs the two exact route certificates and
weights, with a fresh 986797-step literal trajectory to the
second route's bounded section. The ancillary files also compare the
full rescaled cyclic state sets.\\[3pt]
Excluded-length examples
& Substitution in the paired formula at length 23 and a
ten-step affine feeder calculation at length 25.
& Supplementary literal controls in the ancillary files. The arguments
in Section~\ref{sec:remarks} use explicit comparisons.\\[3pt]
Order-8 cohorts
& Exact phase arrays, ray matrices, polytope volumes and
51-digit entrance bounds.
& Independently reads the phase arrays and triangulates
the rational polytopes. It also reconstructs both physical lifts.\\[3pt]
Finite cube evidence
& Replays all sampled classifications and exact binomial
decisions for the order-8 sample and the parity test.
& The finite samples are reproducible. The reconstruction does not infer a
cube limit or reconstruct the sample statistics.\\
\bottomrule
\end{tabular}
\end{center}

\appendix
\section{The five sign words at lengths 15 and 17}\label{app:eight}

In \eqref{eq:root} each \(+\), \(-\), or \(0\) prescribes the sign
of the endpoint \(z\). The negative branch has \(h=-1-z\),
including its admissible boundary \(z=-1\). These are the words
from the primary certificate. An integer in the last column
bounds the spread divided by \(L\) for every permitted copy.
\[
\begin{array}{rrrr}
\text{row} & \text{length} & \text{loss per copy} & \text{spread bound}\\
\midrule
1 & 4 & 0 & 177\\
2 & 16 & 4 & 179\\
3 & 52 & 12 & 155\\
4 & 48 & 18 & 124\\
5 & 52 & 13 & 72
\end{array}
\]
{\small
\begin{verbatim}
++++
+++-+++++++-++++
+++-++++++--+++0++-++++-++0+++-0+++-++-++++-++-++++0
++-++++-++++++-++++-++++++-++++-++++++-++++-++++
++-+0++-++++-++0++-++++-0+++-+-0++0-++-+0++-++-+-++0
\end{verbatim}
}

\begin{lemma}[Finite period-eight certificate]\label{lem:comp8}
At the source of Section~\ref{sec:eight}, the five words with
multiplicities \(a,b,1,c+1,d+1\) satisfy every sign comparison.
Their last state is \((0,0,18,2)\).
The ordinary continuation first becomes periodic after 35 steps
and has least period eight, zero cyclic loss and canonical
relative signature \((0,-1,2,-1,0)\), with the anchor prepended.
The continuation through one terminal period has loss 10 and
spread at most 18.
\end{lemma}

For each word use its unshifted start and the remaining parameter
rays. The primary program verifies that the constant of each
signed affine comparison is at least one and its direction
coefficients are nonnegative, or, for a zero read, that
both constant and coefficients vanish. It checks that the active
rays are restored and the prescribed displacement is realized.
Thus a completed copy leaves the prescribed start for the next
copy. The independent program reconstructs the same comparisons
from the root formula, then checks the floor choices and
profiles in actual gap coordinates.

\section{Finite certificates for the strict lengths}\label{app:strict}

Let \(\operatorname{sp}(x)=\max(0,x)-\min(0,x)\).
In a record \((r,x_*,I,a,b)\), \(I\) lists \(r\) pairs
\((s_i,t_i)\) of insertion endpoints, and \([a,b]\)
specifies a complete terminal return. Write \(x(t)=T^t(x_*)\),
\(H_b=\sum_{0\le t<b}h(x_r(t))\), and
\(M_b=\max_{0\le t\le b}\operatorname{sp}(x(t))\).
The finite checks are the following.

\begin{lemma}[Full-rank reference records]\label{lem:compstrict}
The accompanying records contain one reference orbit for each
\(3\le r\le13\) and a shifted reference at order 12.
Each record satisfies
\begin{enumerate}
\item its insertion intervals are ordered, disjoint, contained in
      \([0,a]\), and have positive lengths divisible by \(r\);
\item \(u_i=S^{-s_i}(x(s_i)-x(t_i))\) is nonnegative, and every label
      with nonzero coefficient in \(u_i\) has positive read at
      every such read before \(t_i\);
\item in first-nonpositive-read order, ray \(i\) vanishes on
      earlier retiring labels and has positive diagonal entry.
      The determinant is nonzero;
\item \(x(a)=x(b)\), with no loss on \([a,b)\).
      The least return periods and the linear loss and spread
      bounds are as follows.
\end{enumerate}
In the shifted record the original slot-7 source increases by
\(12592\). All twelve insertion rays agree with the unshifted record.
\[
\begin{array}{r|rrrr}
r & \text{least period} & b+\sum_i(t_i-s_i) & C_{\rm loss} & C_{\rm spread}\\
\midrule
3 & 2 &2094&311&1924\\
4 & 6 &1222&284&835\\
5 &10 &1782&419&947\\
6 & 5 &26498&9918&10613\\
7 &14 &10820&4019&3901\\
8 &14 &47083&18244&14557\\
9 &18 &58566&22411&17178\\
10&20 &142150&64395&25663\\
11&22 &365572&144048&92554\\
12&24 &1519048&760005&193166\\
13&26 &495253&223003&101757\\
12\text{ shifted}&11&1618064&809022&205758
\end{array}
\]
Here
\[
\begin{aligned}
 C_{\rm loss}
   &=H_b+\sum_i\sum_{s_i\le t<t_i}h(x_r(t)),\\
 C_{\rm spread}
   &=M_b+\sum_i\operatorname{sp}(u_i).
\end{aligned}
\]
Use \(C=\max(C_{\rm loss},C_{\rm spread})\) in \eqref{eq:D}.
\end{lemma}

For the ordinary records \(3\le r\le11\) and \(r=13\)
use \file{root-certificates.jsonl}. The two order-12 files are
\file{certificate.jsonl} and \file{period132-root.jsonl}.
They give \(x_*,I,a,b\) explicitly. In condition 3 the first
nonpositive reads are taken from the unshifted records. For the
shifted record the positivity checks in condition 2 and the
unchanged rays supply the same rank.

\begin{lemma}[Screen and exact lift values]\label{lem:compscreen}
The odd divisors \(n\ge5\) of \(2^r-1\) or \(2^r+1\) of
exact folded order \(r\), for \(3\le r\le13\), give exactly
the 47 rows in Table~\ref{tab:lengths} and the nine exclusions
in Section~\ref{sec:lift}. For a row in the table, the
invariant profile defined below has positive pairing margin
\(\Delta\), positive secondary margin \(\gamma\), and
\(\delta_n=A/2\). Its value of \eqref{eq:D} equals the tabulated \(D_n\).
At \(n=45,91\) the shifted record gives respectively
\(D_{n,132}=92890047744,923554563888\) with the same profiles
and radius thresholds. The complementary permutations have
the orders and orbit lengths in Section~\ref{sec:coexist}.
\end{lemma}

Set \(w_j=0\) on \(J\) and \(w_j=1\) elsewhere, with
the following departures from the default. A representative
is the smallest index of its permutation orbit.
\[
\begin{array}{r|l}
n&\text{representative and assigned weight}\\
\midrule
43 & (3,3),(7,2)\\
45 & (9,2),(15,2)\\
51 & (3,2),(9,3),(17,0)\\
57 & (19,-1)
\end{array}
\]
Put \(\bar w=m^{-1}\sum_j w_j\) and
\(\pi_j=1/m+(w_j-\bar w)/(2m(m+1)W)\).
Here \(W=\max w_j-\min w_j\), except at \(45,91\) where
\(W=\bar w-\min w_j\). This makes
\(\min\pi_j>1/(m+1)\). The verifier checks every retained
witness on all secondary rows. On each excluded forced-tie
length it checks a zero row on every complementary indicator.
Suffix sums and \eqref{eq:D} then require only rational arithmetic.

\section{Certificates in the thirteen-digit base class}\label{app:13}

\begin{lemma}[Six translations to a period-177 phase]\label{lem:comp13a}
In \(X/7\) coordinates the source is
\[
 (2148,646,2,1355,-1340,1922)+\sum_{i=1}^6 a_i u_i .
\]
The file \file{route.json} gives the six \(u_i\) and count increments.
Their determinant is \(84015792\).
The lengths and constant copy counts of the six translating
words are respectively
\[
 (6,12,12,36,36,18)\quad\text{and}\quad(113,7,44,2,36,30).
\]
For all nonnegative integer source parameters the first and last
copy guards and bridge guards hold on the affine pair branches.
Their endpoint is \((0,2,0,1,2,2)\) and all source directions
vanish. On the box \(0\le a_i\le L\), through one terminal
period, the unnormalized radius \(R\) is at most \(22834+5663L\).
At the endpoint the least period is 177 and the maximum \(R\)
on that cycle is 28.
\end{lemma}

\begin{lemma}[Six returns to the bounded section]\label{lem:comp13b}
In \(X\) coordinates define
\[
  S_0=(90594,64043,49868,17220,10073,13923).
\]
The directions \(D_i\) are the columns of
\[
\begin{array}{rrrrrr}
42 & 84 & 210 & 126 & 147 & 84\\
35 & 70 & 28 & 105 & 84 & 21\\
28 & 7 & 91 & 84 & 28 & 7\\
21 & -7 & -91 & -84 & -28 & -7\\
14 & 28 & -77 & 42 & -84 & -21\\
7 & 14 & -112 & 21 & 63 & 63
\end{array}
\]
with determinant of absolute value \(411848913042\).
Set \(B_0=140000\), \(M=63\lfloor(B-B_0)/(63B_0)\rfloor\)
and \(w_i=\rho_i/12\), where
\[
 \rho=(9584/7,169/21,902/21,1192/21,826/9,2063/63).
\]
For \(B\equiv0\pmod7\), \(B\ge8960000\), the box
\((1+M)S_0+\sum_i z_iD_i\), \(|z_i|\le\lfloor Mw_i\rfloor\),
satisfies \(g_0<B/8\). The route has six translating words of
lengths \(6,12,30,18,36,18\), with copy counts respectively
\(1370,8,42,56,89,6\) plus \(M\rho_i+z_i\).
The bridges after the first five have lengths \(0,25,2,100,56\).
The final return is followed by a 16-step partial word.
The route is paired throughout and its first state of \(R\le2000\)
has \(H=(-1988,1484,1491,-1974,315,322,350)\).
That state reaches the same period-177 cycle as in
Lemma~\ref{lem:comp13a} after 437 steps and its whole suffix
has \(R\le1988\).
\end{lemma}

The primary verifier reconstructs each sign pattern from literal
subtraction of the starting state, then checks the branch inequalities
at both ends of each word and the exact affine displacement
and bridge identities. It checks which comparison triggers
each passage, the exterior witness at every prior phase and
the disappearance of all seven slopes at the bounded section.
For the mass in \eqref{eq:mu}, the integral mean of
\(v_i^j\) on \([-w_i,w_i]\) is \(w_i^j/(j+1)\) when \(j\) is even
and zero when \(j\) is odd. Exact expansion gives
\[
\mu =
\frac{\mbox{\scriptsize 3773279180259731755300995131345652558791147811352345556322576361}}
     {\mbox{\scriptsize 3392252763571048539774051887827968}\cdot10^{51}}.
\]
The two terminal cycles agree after converting the
\(X/7\) coordinates of Lemma~\ref{lem:comp13a} to \(X\). At each fixed
base applying \(K_B\) to a phase is a well-defined map back to
words, so equality of the full gap cycle sets gives the same
whole-word cycle.

\section{Exact order-8 components and finite sampling}\label{app:cube}

\begin{lemma}[Cohorts in the order-8 cube]\label{lem:compcube}
The two order-8 records have full-rank nonnegative insertion
matrices \(U_1,U_2\) with determinants of absolute value
\(1450463814294300\) and \(225273027450240\).
The last active labels have descents \(d_1=369\), \(d_2=66\)
on their respective projected returns. Of the \(d_j\)
consecutive incoming heights surviving one return, exactly
\(194,39\), respectively, reach a zero-loss terminal cycle.
The corresponding disjoint rotated cohorts have limiting cube
weights \(8d_j V_j\), where
\[
\begin{split}
V_j&=\operatorname{vol}_8\{a\ge0:U_ja\le\mathbf1\},\\
V_1&=\frac{197}{11088438305193849600000},\\
V_2&=\frac{643703901358114633444077107120707}
 {10584303473669331890415126022076788207622553600000000}.
\end{split}
\]
At \(n=51\) the two guarded entrance paths have radius threshold
\(2/4563\) and values \(D=7997661360,50745451536\).
\end{lemma}

On a projected return let \(a_i\) be the successive active-read
offsets, \(a_*=\min_i a_i\), and \(d>0\) the decrease in the
active label after one turn of the return. The heights
\(1-a_*,\ldots,d-a_*\) all survive that return.
The first failed read in the next return is the first
\(i\) with \(y-d+a_i\le0\). Adding \(d\) shifts the time of
retirement by a complete return and leaves its state unchanged.
Thus one finite enumeration covers all these residues,
and exact repeated-state equality with the cycle-loss sum gives
the terminal classification (zero cyclic loss or not) for each. For each residue the chart is
\(x_*+U_j a+\Delta e_\ell\) with \(a\) a nonnegative integer vector,
\(\ell\) the last active label and \(0\le\Delta<d_j\).
The last column is \(d_j e_\ell\), and triangularity makes these
residue translates disjoint. In coefficient space the cube cut
converges after rescaling to the stated polytope. A rational vertex
enumeration and simplex determinants give the volumes.
All first eight reads are positive on these charts.
Shifting the path by any of those phases preserves the limiting
cube weight and gives eight distinct retirement word orders for
each component. The sixteen orders are distinct. This gives the
factor eight in the limiting cube weights. Consequently
\(8(194V_1+39V_2)\) is a whole-cube lower bound for the zero-loss
liminf, and \(1-8(175V_1+27V_2)\) is a whole-cube upper bound
for its limsup. For each physical lift
use the recorded base-scale box on its individual insertion
certificate, as in Section~\ref{sec:lift}.

\begin{lemma}[Finite sample replay]\label{lem:compsample}
The archived order-8 uniform-cube diagnostic at \(L=10^6\) gives
4305 draws with zero-loss terminal cycles among 8192.
In the separate fixed-size parity experiment of 36 prespecified
comparisons, 21 meet the confirmation threshold, 15
are unresolved and none meet the rejection threshold.
\end{lemma}

The parity data use independent seeded samples of 512 draws in
each cube of two sizes \(20000,200000\) for orders 3 through 20.
The capped run leaves unsolved draws as unknown. Each one-sided
exact binomial threshold uses probability \(1/1440\), with
unknowns charged against the tested direction. The prespecified
comparisons, seeds, classifications and exact decisions are archived in the
ancillary files. These statistics are finite-sample evidence.

\clearpage
\section*{Data availability}
The ancillary files accompanying this paper contain the source programs,
the rational certificates, the output manifests and the replay script used
in Section~\ref{sec:verification}, with the reproduction command. They are
archived on Zenodo as version~1.9
(\href{https://doi.org/10.5281/zenodo.24399274}{doi:10.5281/zenodo.24399274}) and are available in the
\href{https://github.com/galyavin/kaprekar-dynamics/releases/tag/v1.9}{v1.9 release of the Kaprekar dynamics repository on GitHub}.

\section*{Tools}
The research, computations and drafting were carried out with substantial
assistance from AI systems (OpenAI Codex and Anthropic Claude), under the
author's direction and close supervision. The separate verification in
Section~\ref{sec:verification} was carried out in the same way, with programs
written independently of the primary ones. The computations use Python and
C++. The author takes responsibility for the contents.

\begingroup
\raggedright
\begin{thebibliography}{COST}
\bibitem[COST]{COST}
E.~Chen, K.~Ono, R.~E.~Schwartz and D.~S.~Thakur,
\emph{Four-digit Kaprekar dynamics in odd bases},
J. Integer Seq. \textbf{29} (2026), Article 26.4.7.
\href{https://arxiv.org/abs/2606.20439}{arXiv:2606.20439}.

\bibitem[Dahl]{Dahl}
C.~D.~Dahl, \emph{Coarse-grained drift fields and attractor-basin entropy
in Kaprekar's routine}, Entropy \textbf{28} (2026), Article 92.
\href{https://doi.org/10.3390/e28010092}{doi:10.3390/e28010092}.

\bibitem[DZ]{DZ}
P.~Devlin and T.~Zeng, \emph{Maximum distances in the four-digit
Kaprekar process}, Integers \textbf{21} (2021), Paper A97.
\href{https://arxiv.org/abs/2010.11756}{arXiv:2010.11756}.

\bibitem[G1]{G1}
R.~Galyavin, \emph{Cycles of the four-digit Kaprekar map in even bases}, 2026.
\href{https://doi.org/10.5281/zenodo.24358739}{doi:10.5281/zenodo.24358739}.

\bibitem[G2]{G2}
R.~Galyavin, \emph{Arrival times for the four-digit Kaprekar map in even bases}, 2026.
\href{https://doi.org/10.5281/zenodo.24358744}{doi:10.5281/zenodo.24358744}.

\bibitem[G3]{G3}
R.~Galyavin, \emph{A common arrival-time law for five-digit Kaprekar dynamics}, 2026.
\href{https://doi.org/10.5281/zenodo.24358746}{doi:10.5281/zenodo.24358746}.

\bibitem[G4]{G4}
R.~Galyavin, \emph{Majorization in odd-digit Kaprekar dynamics and the nine-digit arrival law}, 2026.
\href{https://doi.org/10.5281/zenodo.24358748}{doi:10.5281/zenodo.24358748}.

\bibitem[G5]{G5}
R.~Galyavin, \emph{Seven-digit Kaprekar dynamics: global law, two clocks, parity}, 2026.
\href{https://doi.org/10.5281/zenodo.24358752}{doi:10.5281/zenodo.24358752}.

\bibitem[G6]{G6}
R.~Galyavin, \emph{Linear-time access for eleven-digit Kaprekar dynamics}, 2026.
\href{https://doi.org/10.5281/zenodo.24358756}{doi:10.5281/zenodo.24358756}.

\bibitem[G7]{G7}
R.~Galyavin, \emph{Arithmetic selection of arrival times in eleven-digit Kaprekar dynamics}, 2026.
\href{https://doi.org/10.5281/zenodo.24371745}{doi:10.5281/zenodo.24371745}.

\bibitem[G8]{G8}
R.~Galyavin, \emph{Eleven-digit Kaprekar arrival laws modulo twelve}, 2026.
\href{https://doi.org/10.5281/zenodo.24396437}{doi:10.5281/zenodo.24396437}.

\bibitem[KD1]{KDO}
A.~Kay and K.~Downes-Ward, \emph{Fixed points and cycles of the Kaprekar
transformation: 1. Odd bases}, J. Integer Seq. \textbf{25} (2022),
Article 22.6.7.
\href{https://cs.uwaterloo.ca/journals/JIS/VOL25/Kay/kay5.html}{cs.uwaterloo.ca/journals/JIS/VOL25/Kay/kay5.html}.

\bibitem[KD2]{KDE}
A.~Kay and K.~Downes-Ward, \emph{Fixed points and cycles of the Kaprekar
transformation: 2. Even bases}, 2024.
\href{https://arxiv.org/abs/2408.12257}{arXiv:2408.12257}.

\bibitem[YM]{YM}
A.~Yamagami and Y.~Matsui, \emph{On 3-adic Kaprekar loops},
JP J. Algebra Number Theory Appl. \textbf{40} (2018), no.~6, 957--1028.
\href{https://doi.org/10.17654/NT040060957}{doi:10.17654/NT040060957}.

\end{thebibliography}
\endgroup
\end{document}
